% Source: https://artofproblemsolving.com/community/c3h1453387_algebra_question
% Stable identity: aops-topic-1453387-q12
Moving a graph $k$ units up adds $k$ to the function; moving it $h$ units right replaces $x$ by $x-h$ (a point $(x_0,y_0)$ moves to $(x_0+h,y_0)$, and $(x_0+h)-h=x_0$).
\begin{enumerate}
\item[1.] $\boxed{y=x^2+3}$.
\item[2.] $\boxed{y=(x+5)^2}$.
\item[3.] $\boxed{y=(x-5)^2-4}$.
\end{enumerate}
In each panel of the figure the dashed grey curve is $y=x^2$ and the coloured curve is the moved graph, with its vertex marked.
\begin{asy}
usepackage("amsmath");
size(11cm);
defaultpen(fontsize(10pt));
import graph;
real f0(real x){return x^2;} real f1(real x){return x^2+3;}
real f2(real x){return (x+5)^2;} real f3(real x){return (x-5)^2-4;}
void panel(pair o, real a, real b, real g(real), real lo, real hi, pen p, pair v, string vl, pair vdir, string cap){
  draw(shift(o)*((a,0)--(b,0)),Arrow(4)); draw(shift(o)*((0,-5)--(0,10)),Arrow(4));
  draw(shift(o)*graph(f0,-sqrt(9),sqrt(9)),gray+dashed);
  draw(shift(o)*graph(g,lo,hi),p+linewidth(1));
  dot(o+v,p); label(vl,o+v+(vdir==E ? (1.7,0) : (0,0)),vdir,p);
  label(cap,o+((a+b)/2,-7.8));
}
panel((0,0),-4,4,f1,-sqrt(6),sqrt(6),blue,(0,3),"$(0,3)$",E,"1. $y=x^2+3$");
panel((16,0),-9,3,f2,-5-sqrt(9),-5+sqrt(9),red,(-5,0),"$(-5,0)$",S,"2. $y=(x+5)^2$");
panel((25,0),-3,10,f3,5-sqrt(13),5+sqrt(13),deepgreen,(5,-4),"$(5,-4)$",S,"3. $y=(x-5)^2-4$");
label("\mbox{}",(min(currentpicture,user=true)+max(currentpicture,user=true))/2,black); // keeps the text after the figure black when LaTeX typesets the labels inline
\end{asy}
